\section{2. Applicable and Reference Documents}\label{applicable-and-reference-documents}

\subsection{2.1. Primary Documents}\label{primary-documents}

\begingroup\GSEGTableStyle{2}
\begin{longtable}{|P{0.43}|P{0.57}|}
\hline\rowcolor{TableHeaderBlue}
\GSEGHeaderText{Source} & \GSEGHeaderText{Role in the DDF}\\ \hline
\endhead
\hline\endlastfoot
EWAS System Report (AD-01) & requirements R-001–R-049, mission, CONOPS, WBS\\ \hline
EWAS-HSB-GSEG-TRS-02, Issue 0.7 & segment requirements and RTM\\ \hline
EWAS-HSB-SSEG-TRS-03 (AD-10) & satellite interfaces and functions, particularly ISL\\ \hline
ECSS-E-ST-10C, ECSS-E-ST-10-06C & systems engineering and requirements documentation\\ \hline
ECSS-E-ST-10-24C Rev.1 (AD-06) & interface management\\ \hline
ECSS-E-ST-40C Rev.1 (AD-07) & software\\ \hline
EWAS System Data Dictionary (RD-06) & product nomenclature\\ \hline
ECSS-Q-ST-30-09C /\allowbreak{} ECSS-Q-ST-40-12C (RD-07/\allowbreak{}RD-08) & availability and fault tree analysis\\*
\end{longtable}
\endgroup

\subsection{2.2. Configuration and ICDs}\label{configuration-and-interface-documentation}

Final protocols, frames, cryptographic boundaries, delays, synchronization, number of channels and acknowledgement rules require agreed ICDs. The DDF source files and registers describe the current assumptions and do not replace the ICDs.

\subsection{2.3. Terminology}\label{terminology}

Names and abbreviations are adopted in accordance with the TRS. In particular, \textbf{GSEG} denotes one geographical site, while \textbf{GIS A/\allowbreak{}B domains} denote two independently maintainable \emph{equipment domains at the same site}; these are not two geographical centers.
